%!TEX TS-program = xelatex
%!TEX encoding = UTF-8 Unicode
% Awesome CV LaTeX Template for CV/Resume
%
% This template has been downloaded from:
% https://github.com/posquit0/Awesome-CV
%
% Author:
% Claud D. Park <posquit0.bj@gmail.com>
% http://www.posquit0.com
%
% Template license:
% CC BY-SA 4.0 (https://creativecommons.org/licenses/by-sa/4.0/)
%


%-------------------------------------------------------------------------------
% CONFIGURATIONS
%-------------------------------------------------------------------------------
% A4 paper size by default, use 'letterpaper' for US letter
\documentclass[11pt, letter]{awesome-cv}

% xdvipdfmx embeds this template's CID-keyed OpenType text (via fontspec) without
% any actual space glyph in the content stream — word gaps are pure TJ-array
% position adjustments (confirmed: the font's CID map has zero CIDs mapped to
% U+0020), and paragraph-mode text keeps each word as its own glyph run even with
% \XeTeXinterwordspaceshaping=2 (tried; only merges runs inside single-line boxes
% that never had a line-break opportunity, e.g. entry titles, not multi-line
% bullets). This appears to be inherent to how XeLaTeX ships out justified/ragged
% paragraphs with embedded OpenType fonts, not something fixable from the .tex
% side. \XeTeXgenerateactualtext=1 is kept since it's a harmless, real
% improvement for copy/paste and in-PDF search fidelity.
%
% The actual fix for text-extraction robustness lives in
% scripts/validate_ats.py: it now reads pdftotext's default (layout-reconstructed)
% output instead of -raw, which infers word boundaries from glyph geometry across
% the whole line rather than trusting the size of a single TJ-array number, and
% agrees with pdfminer.six's independent geometric reconstruction. See that file
% for details.
\XeTeXgenerateactualtext=1

% Configure page margins with geometry
\geometry{left=1.4cm, top=.6cm, right=1.4cm, bottom=1.0cm, footskip=.5cm}

% Color for highlights
% Awesome Colors: awesome-emerald, awesome-skyblue, awesome-red, awesome-pink, awesome-orange
%                 awesome-nephritis, awesome-concrete, awesome-darknight
\colorlet{awesome}{awesome-orange}
% Uncomment if you would like to specify your own color
% \definecolor{awesome}{HTML}{3E6D9C}

% Colors for text
% Uncomment if you would like to specify your own color
% \definecolor{darktext}{HTML}{414141}
% \definecolor{text}{HTML}{333333}
% \definecolor{graytext}{HTML}{5D5D5D}
% \definecolor{lighttext}{HTML}{999999}
% \definecolor{sectiondivider}{HTML}{5D5D5D}

% Set false if you don't want to highlight section with awesome color
\setbool{acvSectionColorHighlight}{true}

% If you would like to change the social information separator from a pipe (|) to something else
\renewcommand{\acvHeaderSocialSep}{\quad\textbar\quad}


%-------------------------------------------------------------------------------
%	PERSONAL INFORMATION
%	Comment any of the lines below if they are not required
%-------------------------------------------------------------------------------
% Available options: circle|rectangle,edge/noedge,left/right
% \photo[rectangle,edge,right]{./examples/profile}

% Real contact info (phone number, non-placeholder email) for job-application builds
% lives in resume/personal-info.tex, which is gitignored and never committed to this
% public repo. If that file isn't present (e.g. in CI, or a fresh clone), these
% placeholder defaults are used instead.
\newcommand{\resumemobile}{425-555-0100}
\newcommand{\resumeemail}{j43cheun@uwaterloo.ca}
\IfFileExists{resume/personal-info.tex}{\input{resume/personal-info.tex}}{}

\name{Justin}{Cheung}
\position{Senior Software Engineer}
\address{Redmond, WA, USA}

\mobile{\resumemobile}
\email{\resumeemail}
%\dateofbirth{January 1st, 1970}
%\homepage{j43cheun.github.io}
\github{j43cheun}
\linkedin{j43cheun}
% \gitlab{gitlab-id}
% \stackoverflow{SO-id}{SO-name}
% \twitter{@twit}
% \skype{skype-id}
% \reddit{reddit-id}
% \medium{madium-id}
% \kaggle{kaggle-id}
% \hackerrank{hackerrank-id}
% \googlescholar{googlescholar-id}{name-to-display}
%% \firstname and \lastname will be used
% \googlescholar{googlescholar-id}{}
% \extrainfo{extra information}

%\quote{``Be the change that you want to see in the world."}

%-------------------------------------------------------------------------------
\begin{document}

% Print the header with above personal information
% Give optional argument to change alignment(C: center, L: left, R: right)
\makecvheader[C]

% Print the footer with 3 arguments(<left>, <center>, <right>)
% Leave any of these blank if they are not needed
%\makecvfooter
%  {\today}
%  {Byungjin Park~~~·~~~Résumé}
%  {\thepage}

\vspace{\acvSectionContentTopSkip}
\vspace{-4.5mm}

%-------------------------------------------------------------------------------
%	CV/RESUME CONTENT
%	Each section is imported separately, open each file in turn to modify content
%-------------------------------------------------------------------------------
\input{resume/technical-skills.tex}

\vspace{\acvSectionContentTopSkip}
\vspace{-4.5mm}

%-------------------------------------------------------------------------------
%	SECTION TITLE
%-------------------------------------------------------------------------------
\cvsection{Experience}


%-------------------------------------------------------------------------------
%	CONTENT
%-------------------------------------------------------------------------------
\begin{cventries}

  \cventrytwopositions
  {Software Engineer II - Storage Spaces Dev Team} % Job title 1
  {Microsoft Corporation} % Organization
  {Redmond, WA, USA} % Location
  {Sep. 2015 - Jan. 2022} % Date(s) for Job title 1
  {
    \begin{cvitems} % Description(s) of tasks/responsibilities
      \item {
        Optimized Storage Spaces Repair (i.e., RAID resilvering) by enhancing Storage
        Spaces metadata to track stale logical data copies at stripe granularity
        (i.e., 256 KiB) vs. extent granularity (i.e., 1 GiB) previously, which
        reduced its run time from hours to minutes for multi-terabyte Storage Spaces.
      }
      \item {
        Implemented Storage Spaces support for node-mirrored mirror and node-mirrored
        parity multilevel resiliency on 2-node Storage Spaces Direct (S2D) Failover
        Clusters, enabling them to tolerate 1 node + 1 drive failure versus 1 node
        failure previously.
      }
      \item {
        Designed and implemented an extra-resilient cache for staging writes to
        Storage Space ranges backed by a suboptimal number of logical copies to
        enable write forward progress without compromising on redundancy and risking
        data loss.
      }
    \end{cvitems}
  }
  {Senior Software Engineer - Hyper-V Modern Virtualization Team} % Job title 2
  {Jan. 2022 - Present} % Date(s) for Job title 2
  {
    \begin{cvitems} % Description(s) of tasks/responsibilities
      \item {
        Implemented Hyper-V virtualization stack support for Virtualization
        Infrastructure Driver (VID.sys) runtime reload during VM Preserving Host
        Update (VM-PHU) Ultralite with Handle Brokering, which enables VID.sys to be
        runtime reloaded during VM-PHU Ultralite without draining I/O.
      }
      \item {
        Contributed to the design and implementation of Virtual Processor (VP) Auto
        Suspend for VM-PHU, which defers VP explicit suspend to the Microsoft Hypervisor until
        the arrival of root partition intercepts, thereby minimizing VM compute
        blackout during host servicing and enabling Hyper-V virtualization stack
        components to be serviced on the host while VPs are still running.
      }
      \item {
        Provisioned Azure Host OS media with slipstreamed software updates, enabling
        inner-loop test coverage for VM-PHU Self upgrades from down-level to
        active-branch builds and, more broadly, validation of changes against Azure's
        actual OS composition locally — previously only possible in Azure itself due
        to the complexity of replicating that composition on lab machines.
      }
    \end{cvitems}
  }

  \cventry
    {Compiler Optimization Developer (Intern)} % Job title
    {IBM Canada Ltd.} % Organization
    {Markham, ON, Canada} % Location
    {Sep. 2014 - Dec. 2014} % Date(s)
    {
      \begin{cvitems} % Description(s) of tasks/responsibilities
        \item {
          Implemented portions of a translator from LLVM IR to Wcode, IBM's internal
          compiler intermediate representation, to evaluate whether an LLVM-based
          frontend could target IBM's XL C/C++/Fortran compiler's backend and
          optimizer.
        }
        \item {
          Expanded test coverage for the LLVM IR translator by writing a C/C++ shim
          that emulated the test drivers of IBM's XL C/C++/Fortran compiler, allowing
          existing test buckets to be reused for validation.
        }
      \end{cvitems}
    }

  \cventry
    {System Software Developer (Intern)} % Job title
    {Nav Canada} % Organization
    {Ottawa, ON, Canada} % Location
    {Apr. 2012 - Aug. 2012} % Date(s)
    {
      \begin{cvitems} % Description(s) of tasks/responsibilities
        \item {
          Ported Nav Canada's build distribution tool for air traffic control software
          from HP-UX to Red Hat Enterprise Linux with a Java ICEFaces web frontend,
          replacing a complex, largely manual process with an automated one.
        }
      \end{cvitems}
    }

\end{cventries}


\vspace{\acvSectionContentTopSkip}
\vspace{-4.5mm}

%\input{resume/honors.tex}
%\input{resume/certificates.tex}
% \input{resume/presentation.tex}
% \input{resume/writing.tex}
% \input{resume/committees.tex}

%-------------------------------------------------------------------------------
%	SECTION TITLE
%-------------------------------------------------------------------------------
\cvsection{Patents}


%-------------------------------------------------------------------------------
%	CONTENT
%-------------------------------------------------------------------------------
\begin{cventries}

  \cventry
    {Kevin Michael Broas, Sai Ganesh Ramachandran, Vignesh Kudigrama Shenoy, Justin Sing Tong Cheung} % Institution
    {Virtual Processor Auto-Suspend During Virtualization Stack Servicing}
    {}
    {Apr. 24, 2024} % Date(s)
    {
      \begin{cvitems}
        \item {
          Application Number: US 18/645,012. Publication Number: US 2025/0335227 (patent pending).
        }
      \end{cvitems}
    }

  \cventry
    {Taylor Alan Hope, Vinod R. Shankar, Justin Sing Tong Cheung} % Institution
    {Extra-Resilient Cache for Resilient Storage Array}
    {}
    {Jun. 28, 2022} % Date(s)
    {
      \begin{cvitems}
        \item {
          Application Number: US 17/100,557. Patent number: US 11372557.
        }
      \end{cvitems}
    }

  \cventry
    {Karan Mehra, Justin Sing Tong Cheung, Vinod R. Shankar, Grigory Borisovich Lyakhovitskiy} % Institution
    {Multilevel Resiliency}
    {}
    {Oct. 20, 2020} % Date(s)
    {
      \begin{cvitems}
        \item {
          Application Number: US 16/389,605. Patent number: US 10809940.
        }
      \end{cvitems}
    }

\end{cventries}

\vspace{\acvSectionContentTopSkip}
\vspace{-4.5mm}

\input{resume/education.tex}

\vspace{\acvSectionContentTopSkip}
\vspace{-4.5mm}

% \input{resume/extracurricular.tex}


%-------------------------------------------------------------------------------
\end{document}